\chapter{Supermartingales from independent and conditionally sub-Gaussian increments}
\label{chap:pre_martingale}

Throughout, $(\Omega, \mathcal A, \Prob)$ is a probability space with a filtration $\F$ indexed by
$\N$. A real process $Y$ has \emph{independent increments along $\F$} if each $Y_n$ is
$\F_{n+1}$-measurable and independent of $\F_n$. When $\Omega$ is a standard Borel space, a
sub-$\sigma$-algebra $\mathcal G$ has a conditional expectation kernel $\kappa_{\mathcal G}$ (a
regular version of $\Prob(\cdot \mid \mathcal G)$, Mathlib's \texttt{condExpKernel}), and a random
variable $Y$ is \emph{conditionally $c$-sub-Gaussian given $\mathcal G$} if $e^{tY}$ is integrable
for all $t$ and, for almost every $\omega$, $\int e^{tY} \, d\kappa_{\mathcal G}(\omega) \le
e^{ct^2/2}$ for all $t \in \R$ (Mathlib's \texttt{HasCondSubgaussianMGF}).

\section{Freezing and supermartingales}

\begin{lemma}[Freezing lemma]
  \label{lem:lintegral_comp_indep}
  \lean{ProbabilityTheory.setLIntegral_comp_of_indep, ProbabilityTheory.Indep.comap_comp}
  \leanok
  Let $Y$ be a random variable independent of a sub-$\sigma$-algebra $\mathcal G$, $Z$ a
  $\mathcal G$-measurable random variable, $F \ge 0$ measurable and $A \in \mathcal G$. Then
  $\int_A F(Z, Y) \, d\Prob = \int_A \int F(Z, y) \, d\Prob_Y(y) \, d\Prob$. Moreover, every
  measurable function of $Y$ is independent of $\mathcal G$.
\end{lemma}

\begin{proof}
  \leanok
  $(Z, \indic_A)$ is $\mathcal G$-measurable, hence independent of $Y$, so the law of
  $((Z, \indic_A), Y)$ is the product of the marginals. Apply Tonelli's theorem to
  $((z, a), y) \mapsto a F(z, y)$.
\end{proof}

\begin{lemma}[Freezing lemma, conditional form]
  \label{lem:lintegral_comp_condExpKernel}
  \lean{ProbabilityTheory.lintegral_comp_condExpKernel}
  \leanok
  Let $\Omega$ be standard Borel, $\mathcal G$ a sub-$\sigma$-algebra, $Z$ a $\mathcal G$-measurable
  random variable and $F \ge 0$ measurable. Then
  $\int F(Z(\omega), \omega) \, d\Prob(\omega) = \int \int F(Z(\omega), \omega') \,
  d\kappa_{\mathcal G}(\omega)(\omega') \, d\Prob(\omega)$.
\end{lemma}

\begin{proof}
  \leanok
  The composition-product of $\Prob|_{\mathcal G}$ and $\kappa_{\mathcal G}$ is the law of
  $\omega \mapsto (\omega, \omega)$ on $(\Omega, \mathcal G) \times (\Omega, \mathcal A)$ (Mathlib's
  \texttt{compProd\_trim\_condExpKernel}). Apply Tonelli's theorem for composition-products to
  the $\mathcal G \otimes \mathcal A$-measurable function $(\omega, \omega') \mapsto F(Z(\omega),
  \omega')$; the inner integral is $\mathcal G$-measurable, so its integral against
  $\Prob|_{\mathcal G}$ is its integral against $\Prob$.
\end{proof}

\begin{lemma}[Supermartingales with independent increments]
  \label{lem:supermartingale_of_markov}
  \lean{MeasureTheory.supermartingale_of_setLIntegral_succ_le,
    MeasureTheory.supermartingale_of_lintegral_le}
  \leanok
  \uses{lem:lintegral_comp_indep}
  An adapted, a.s.\ nonnegative process $V$ with $V_0$ integrable and $\int_A V_{n+1} \le \int_A
  V_n$ for all $n$ and $A \in \F_n$ is a supermartingale. In particular, let $V$ be such a process
  (without the inequality), $Y$ with independent increments along $\F$, and for each $n$ an
  $\F_n$-measurable $Z_n$ and a measurable $\Phi_n \ge 0$ with $V_{n+1} = \Phi_n(Z_n, Y_n)$ a.s.\
  and $\int \Phi_n(Z_n, y) \, d\Prob_{Y_n}(y) \le V_n$ a.s. Then $V$ is a supermartingale.
\end{lemma}

\begin{proof}
  \leanok
  With $A = \Omega$ and induction, every $V_n$ is integrable. Conclude with the characterization of
  supermartingales by set integrals (Mathlib's \texttt{supermartingale\_of\_setIntegral\_succ\_le}).
  For the second statement, for $A \in \F_n$, Lemma~\ref{lem:lintegral_comp_indep} gives
  $\int_A V_{n+1} = \int_A \int \Phi_n(Z_n, y)\,d\Prob_{Y_n}(y) \le \int_A V_n$.
\end{proof}

\begin{lemma}[Convergence of nonnegative supermartingales]
  \label{lem:supermartingale_tendsto}
  \lean{MeasureTheory.Supermartingale.exists_ae_tendsto_of_nonneg}
  \leanok
  A supermartingale $V$ with $V_n \ge 0$ a.s.\ converges almost surely to a finite limit
  $V_\infty$, and $\E[V_\infty] \le \E[V_0]$.
\end{lemma}

\begin{proof}
  \leanok
  $-V$ is a submartingale with $\E\abs{V_n} = \E V_n \le \E V_0$; apply Mathlib's martingale
  convergence theorem (\texttt{Submartingale.ae\_tendsto\_limitProcess}). Fatou's lemma gives the
  bound on the expectation.
\end{proof}

\begin{lemma}[Mixtures of supermartingales]
  \label{lem:supermartingale_integral}
  \lean{MeasureTheory.supermartingale_integral_of_integrable, MeasureTheory.supermartingale_integral}
  \leanok
  Let $\pi$ be an s-finite measure on a measurable space $L$ and $(f_\ell)_{\ell \in L}$ a family of
  nonnegative supermartingales such that $(\ell, \omega) \mapsto f_\ell(n, \omega)$ is
  $\mathcal B(L) \otimes \F_n$-measurable for each $n$ and $\ell \mapsto \E[f_\ell(0)]$ is
  $\pi$-integrable (for instance $\pi$ finite and $\E[f_\ell(0)] \le C$ for all $\ell$). Then
  $n \mapsto \int f_\ell(n, \cdot)\, d\pi(\ell)$ is a supermartingale.
\end{lemma}

\begin{proof}
  \leanok
  For $A \in \F_n$, by Tonelli's theorem,
  $\int_A \int f_\ell(n+1) \, d\pi \, d\Prob = \int \int_A f_\ell(n+1) \, d\Prob \, d\pi \le \int
  \int_A f_\ell(n) \, d\Prob \, d\pi = \int_A \int f_\ell(n) \, d\pi \, d\Prob$, and
  $\E[f_\ell(n)] \le \E[f_\ell(0)]$ gives integrability.
\end{proof}

\section{Conditionally sub-Gaussian sums}

In this section each $Y_n$ is $\F_{n+1}$-measurable, $w$ is an adapted real process,
$T_n := \sum_{k<n} w_k Y_k$ and $A_n := \sum_{k<n} w_k^2$, and the increments satisfy the
\emph{moment bound at predictable parameters}: for every $n$, every $\F_n$-measurable $s$ and every
$\F_n$-measurable $g \ge 0$,
\begin{equation}
  \label{eq:moment_bound}
  \E[g \, e^{s Y_n}] \le \E[g \, e^{c s^2 / 2}].
\end{equation}
By the next lemma, this holds when $Y$ has independent $c$-sub-Gaussian increments along $\F$, on
any probability space, and when each $Y_n$ is conditionally $c$-sub-Gaussian given $\F_n$ (a
sub-Gaussian martingale difference sequence), on a standard Borel space. The results of this
section are stated in Lean in these three forms: under~\eqref{eq:moment_bound} (suffix
\texttt{\_of\_lintegral\_le}), for conditionally sub-Gaussian increments (no suffix) and for
independent increments (suffix \texttt{\_of\_indep}).

\begin{lemma}[Moment bound at predictable parameters]
  \label{lem:subg_moment_bound}
  \lean{ProbabilityTheory.HasCondSubgaussianMGF.lintegral_mul_exp_le,
    ProbabilityTheory.HasSubgaussianMGF.lintegral_mul_exp_le_of_indep,
    ProbabilityTheory.HasSubgaussianMGF.hasCondSubgaussianMGF_of_indep,
    ProbabilityTheory.lintegral_mul_exp_le_of_hasCondSubgaussianMGF,
    ProbabilityTheory.lintegral_mul_exp_le_of_indep}
  \leanok
  \uses{lem:lintegral_comp_indep, lem:lintegral_comp_condExpKernel}
  Let $\mathcal G$ be a sub-$\sigma$-algebra and $Y$ a measurable random variable, which is either
  $c$-sub-Gaussian and independent of $\mathcal G$, or conditionally $c$-sub-Gaussian given
  $\mathcal G$ ($\Omega$ standard Borel). Then for all $\mathcal G$-measurable $s$ and $g \ge 0$,
  $\E[g \, e^{sY}] \le \E[g \, e^{cs^2/2}]$. On a standard Borel space, a $c$-sub-Gaussian random
  variable independent of $\mathcal G$ is conditionally $c$-sub-Gaussian given $\mathcal G$.
\end{lemma}

\begin{proof}
  \leanok
  Freeze $(g, s)$: by Lemma~\ref{lem:lintegral_comp_indep} (with $A = \Omega$), resp.\
  Lemma~\ref{lem:lintegral_comp_condExpKernel}, $\E[g \, e^{sY}]$ is the expectation of
  $g(\omega) \int e^{s(\omega) y} \, d\Prob_Y(y)$, resp.\ of $g(\omega) \int e^{s(\omega) Y} \,
  d\kappa_{\mathcal G}(\omega)$, and the inner integral is at most $e^{c s(\omega)^2/2}$ (for almost
  every $\omega$ in the conditional case, since the bound holds for all $t$ simultaneously). For the
  last claim: for rational $q$, almost surely $\int e^{qY} \, d\kappa_{\mathcal G} = \E[e^{qY} \mid
  \mathcal G] = \E[e^{qY}] \le e^{cq^2/2}$ by independence; the moment generating function under
  $\kappa_{\mathcal G}(\omega)$ is continuous, which gives the bound for all real $t$ (Mathlib's
  \texttt{Kernel.HasSubgaussianMGF.of\_rat}).
\end{proof}

\begin{lemma}[Exponential supermartingale]
  \label{lem:exp_subg_supermartingale}
  \lean{ProbabilityTheory.supermartingale_exp_sum_mul_of_lintegral_le,
    ProbabilityTheory.supermartingale_exp_sum_mul,
    ProbabilityTheory.supermartingale_exp_sum_mul_of_indep}
  \leanok
  \uses{lem:supermartingale_of_markov, lem:subg_moment_bound}
  For every $t \in \R$, $U^t_n := \exp(t T_n - c t^2 A_n / 2)$ is a nonnegative supermartingale.
\end{lemma}

\begin{proof}
  \leanok
  $U^t_{n+1} = U^t_n e^{-c s^2/2} e^{s Y_n}$ with the $\F_n$-measurable parameter $s = t w_n$. For
  $A \in \F_n$, \eqref{eq:moment_bound} with $g = \indic_A U^t_n e^{-cs^2/2}$ gives $\int_A U^t_{n+1}
  \le \int_A U^t_n$. Conclude with the first part of Lemma~\ref{lem:supermartingale_of_markov}.
\end{proof}

\begin{lemma}[Convergence on finite quadratic variation]
  \label{lem:subg_sum_tendsto}
  \lean{ProbabilityTheory.ae_exists_tendsto_sum_mul_of_lintegral_le,
    ProbabilityTheory.ae_exists_tendsto_exp_sum_mul_of_lintegral_le,
    ProbabilityTheory.ae_exists_tendsto_sum_mul, ProbabilityTheory.ae_exists_tendsto_sum_mul_of_indep}
  \leanok
  \uses{lem:exp_subg_supermartingale, lem:supermartingale_tendsto}
  Almost surely, if $\sum_n w_n^2 < \infty$ then $T_n$ converges.
\end{lemma}

\begin{proof}
  \leanok
  By Lemmas~\ref{lem:exp_subg_supermartingale} and~\ref{lem:supermartingale_tendsto}, $U^1$ and
  $U^{-1}$ converge a.s.\ to finite limits $L_{\pm}$. On $\{A_\infty < \infty\}$:
  $e^{T_n} = U^1_n e^{cA_n/2} \to L_+ e^{cA_\infty/2}$, and $e^{-T_n} = U^{-1}_n e^{cA_n/2}$ is
  bounded, so $T_n$ is bounded below and the limit of $e^{T_n}$ is positive; take logarithms.
\end{proof}

\begin{lemma}[Little-$o$ on infinite quadratic variation]
  \label{lem:subg_sum_little_o}
  \lean{ProbabilityTheory.ae_forall_abs_sum_mul_le_of_lintegral_le,
    ProbabilityTheory.ae_forall_abs_sum_mul_le, ProbabilityTheory.ae_forall_abs_sum_mul_le_of_indep}
  \leanok
  \uses{lem:exp_subg_supermartingale, lem:supermartingale_tendsto}
  Almost surely, for every $\eps > 0$ there is $C$ with $\abs{T_n} \le C + \eps A_n$ for all $n$.
\end{lemma}

\begin{proof}
  \leanok
  For $t > 0$, $U^{\pm t}$ converges a.s., hence is bounded by some $K_t$, so
  $\pm T_n \le \frac{\log K_t}{t} + \frac{ct}{2} A_n$. Apply this to the countably many
  $t_j = \frac{1}{j+1}$ and, given $\eps$, choose $j$ with $\frac{c}{2(j+1)} \le \eps$; conclude with
  $A_n \ge 0$.
\end{proof}

\section{Hoeffding's inequality and almost sure rates}

\begin{lemma}[Hoeffding's lemma]
  \label{lem:hoeffding_lemma_null}
  \lean{ProbabilityTheory.hasSubgaussianMGF_of_mem_Icc_of_integral_eq_zero}
  \mathlibok
  A centered random variable with values in $[a, b]$ is $\frac{(b-a)^2}{4}$-sub-Gaussian.
\end{lemma}

\begin{lemma}[Almost sure rate of i.i.d.\ sub-Gaussian sums]
  \label{lem:hoeffding_rate}
  \lean{ProbabilityTheory.ae_eventually_abs_sum_range_lt,
    ProbabilityTheory.ae_isBigO_sum_range_sqrt_mul_log}
  \leanok
  Let $Y_n$ be independent and $c$-sub-Gaussian, and $b_n \ge 0$ for large $n$ with
  $\sum_n \exp(-b_n^2/(2cn)) < \infty$. Then almost surely $\abs{\sum_{k<n} Y_k} < b_n$ for all
  large $n$. In particular, almost surely $\sum_{k<n} Y_k = O(\sqrt n \log n)$.
\end{lemma}

\begin{proof}
  \leanok
  Each $Y_n$ is also $(c+1)$-sub-Gaussian, so we may assume $c > 0$.
  By Hoeffding's inequality (Mathlib's \texttt{measure\_sum\_range\_ge\_le\_of\_iIndepFun}), applied
  to $Y$ and $-Y$, $\Prob(\abs{S_n} \ge \sqrt n \log n) \le 2 \exp(-\log^2 n/(2c)) \le 2n^{-2}$ for
  $n$ large. By the Borel--Cantelli lemma, almost surely $\abs{S_n} < \sqrt n \log n$ for all large
  $n$.
\end{proof}
